\chapter{Lua API-referencia}
\section{Futási modell}
Az app.lua egy M táblát ad vissza. Az M.loop() kötelező és ismételten lefut; az M.init(), M.event(event\_type) és M.cleanup() opcionális. loop nélkül a firmware elutasítja az alkalmazást. Futási hiba esetén tartalék alkalmazásra válthat.

A sensor/*.lua érzékelőszkript sensor táblát ad vissza, amelynek sensor.get(index) függvénye számértéket szolgáltat. A firmware a sensor táblát globálisan elérhetővé teszi. Az apppath az aktuális alkalmazás mappájának útvonala. A modulok require("module") hívással tölthetők be, opcionális modulnál használható a pcall(require, "module").
\section{Erőforrások és megjelenítés}
A loadimg és loadfont azonosítót vagy nil értéket ad vissza; a free felszabadítja a betöltött képet vagy betűkészletet. A fill a teljes képernyőt kitölti. Az update továbbítja a képkockát, majd feketére törli a képernyőt. A drawpolygon páros számú koordinátát igényel, legalább négy elemmel. Képkockánként egyszer hívja az update függvényt, az erőforrásokat a cleanup szabadítsa fel.
\section{További modulok}
A system kezeli a késleltetést, háttérvilágítást, V2 felhúzást, 5 V-os kimenetet és alvó módot, valamint olvassa a gombokat, ADC-t és időt. Az onewire DS18B20 érzékelőket keres, és Celsius-fokban olvassa a hőmérsékletet. Az nvs az app NVS-névteret használja; a get szöveget vagy nil értéket, a push logikai értéket ad vissza. A ble inicializálja a BLE-t, és fordulatszámot, sebességet, fokozatot szolgáltat.
\section{display}
\begin{itemize}
\item display.loadimg(path) -> img\_handle|nil
\item display.loadfont(path) -> font\_handle|nil
\item display.free(handle)
\item display.fill(color)
\item display.update()
\item display.drawimg(img, x, y)
\item display.string(font, x, y, text, color)
\item display.fillrect(x, y, width, height, color)
\item display.drawline(x1, y1, x2, y2, color)
\item display.drawcircle(x, y, r, color)
\item display.drawfillcircle(x, y, r, color)
\item display.drawroundrect(x1, y1, x2, y2, r, color)
\item display.drawarrow(x0, y0, x1, y1, w, color)
\item display.drawfillarrow(x0, y0, x1, y1, w, color)
\item display.drawpolygon(\{x1, y1, x2, y2, ...\}, color)
\item display.RED
\item display.GREEN
\item display.BLUE
\item display.WHITE
\item display.BLACK
\item display.YELLOW
\item display.CYAN
\item display.PURPLE
\item display.GRAY
\end{itemize}
\section{system}
\begin{itemize}
\item system.delay(ms)
\item system.backlight(level)
\item system.V2\_pullup(enable)
\item system.out5V(enable)
\item system.keys() -> int
\item system.adc\_ch(channel) -> number
\item system.get\_time() -> hour, minute, second
\item system.config\_sleep(enable)
\end{itemize}
\section{onewire}
\begin{itemize}
\item onewire.scan18b20() -> count
\item onewire.read18b20(index) -> temperature\_c
\end{itemize}
\section{nvs}
\begin{itemize}
\item nvs.get(key) -> string|nil
\item nvs.push(key, value\_string) -> boolean
\end{itemize}
\section{ble}
\begin{itemize}
\item ble.init() -> true
\item ble.get\_telemetry() -> rpm, speed, gear
\end{itemize}
\section{Példák és tanácsok}
A voltmérőpélda a 160\,\texttimes\,80 képpontos területen belüli koordinátákat használ, és korlátozza a sáv szélességét. Az érzékelőpélda ADC-olvasást és simítást mutat, nem teljes NTC-hőmérséklet-számítást: a mintabeli átszámítást valódi kalibrációra kell cserélni. A sensor.get legyen gyors; a kalibráció nvs-ben tárolható, a szöveges értékeket szükség esetén számmá kell alakítani. Az útvonalakat apppath alapján képezze. A DS18B20 keresést egyszer végezze el, majd a ciklusban olvassa az érzékelőket.
\begin{verbatim}
local M = {}
local font = display.loadfont(apppath .. "/../fonts/H24XB.FNT")

local function read_voltage()
return sensor.get(1)
end

function M.loop()
local v = read_voltage()
display.fill(display.BLACK)
display.string(font, 10, 30,
    string.format("%2.1f V", v), display.WHITE)
display.fillrect(10, 60,
    math.max(0, math.min(140, math.floor((v - 9.5) * 12))),
    8, display.GREEN)
display.update()
end

function M.cleanup()
display.free(font)
end

return M
\end{verbatim}
\begin{verbatim}
sensor = {}

local filt = 0.0

function sensor.get(i)
i = i or 0
system.V2_pullup(1)

local raw = system.adc_ch(i)
filt = filt + (raw - filt) / 10.0

local value = filt * 1.0
return value
end

return sensor
\end{verbatim}
